\chapter*{ВВЕДЕНИЕ}
\addcontentsline{toc}{chapter}{ВВЕДЕНИЕ}

Современный автомобильный транспорт стремительно развивается в сторону повышения энергетической эффективности и снижения воздействия на окружающую среду. Одним из наиболее перспективных направлений этого процесса является развитие гибридных силовых установок, объединяющих двигатель внутреннего сгорания и электропривод. Такая синергия различных источников энергии открывает широкие возможности для тонкой настройки режимов работы, повышения общей отдачи системы и более рационального использования топливно-энергетических ресурсов.

Гибридные технологии отличаются большим разнообразием архитектурных решений и концептуальных подходов. Существуют мягкие гибридные системы, обеспечивающие поддержку традиционного двигателя, полноформенные гибриды, способные кратковременно работать исключительно на электрической тяге, а также подключаемые варианты, ориентированные на расширение электрического пробега. С другой стороны, встречаются разные концепции совместной работы традиционного и электрического привода: последовательный, параллельный или комбинированный. Это многообразие конструкций формирует богатый спектр инженерных стратегий, позволяющих адаптировать силовую установку под различные эксплуатационные сценарии и характер движения.

Особое значение приобретает анализ применяемости тех или иных гибридных решений в реальных условиях — от плотного городского трафика до протяжённых загородных маршрутов. Именно в зависимости от режима движения, профиля маршрута и стиля вождения раскрываются сильные стороны каждой концепции: эффективность рекуперации энергии, снижение пиковых нагрузок на двигатель, оптимизация расхода топлива и улучшение динамических характеристик.

Таким образом, целью данной работы становится определение и классификация семейств гибридных установок, поиск оптимального соответствия между данными семействами и условиями эксплуатации транспортных средств.